\documentclass{article}
\usepackage{amsmath,amssymb,amsthm}
\usepackage{algorithm}
\usepackage{algorithmic}
\usepackage{enumitem}
\usepackage{hyperref}
\usepackage{geometry}
\geometry{margin=1in}
\newtheorem{theorem}{Theorem}
\newtheorem{lemma}{Lemma}
\newtheorem{assumption}{Assumption}
\newtheorem{corollary}{Corollary}
\begin{document}

\section*{Random Block Coordinate Descent (RBCD) for Non‑Convex Smooth Objectives}

\section*{Mathematical Formulation}
Let $f:\mathbb{R}^{d}\rightarrow\mathbb{R}$ be differentiable (possibly non‑convex).  
Partition the variable $w\in\mathbb{R}^{d}$ into $B$ disjoint blocks 
\[
w=\bigl(w^{(1)},w^{(2)},\dots ,w^{(B)}\bigr),\qquad 
w^{(i)}\in\mathbb{R}^{d_i},\;\; \sum_{i=1}^{B}d_i=d .
\]
For block $i$ denote the partial gradient by $\nabla_i f(w)\in\mathbb{R}^{d_i}$.

At iteration $t\ge 0$ a block index $i_t\in\{1,\dots ,B\}$ is drawn independently
from a probability distribution $P\{i_t=i\}=p_i>0$, $\sum_i p_i=1$ (uniform sampling corresponds to $p_i=1/B$).

The update rule is
\[
\boxed{
w_{t+1}^{(i_t)} \;=\; w_{t}^{(i_t)} \;-\; \eta_t \,\nabla_{i_t} f(w_t),\qquad 
w_{t+1}^{(j)} = w_{t}^{(j)}\;\;\forall j\neq i_t,
}
\tag{1}
\]
where $\{\eta_t\}_{t\ge0}$ is a stepsize (learning‑rate) sequence, typically constant
$\eta_t\equiv \eta>0$.

\section*{Pseudocode}
\begin{algorithm}[H]
\caption{Random Block Coordinate Descent (RBCD)}\label{alg:RBCD}
\begin{algorithmic}[1]
\STATE \textbf{Input:} initial point $w_0$, stepsize $\eta>0$, block probabilities $\{p_i\}_{i=1}^{B}$, max iterations $T$.
\FOR{$t=0,\dots ,T-1$}
    \STATE Sample block $i_t\sim\{1,\dots ,B\}$ with $\Pr(i_t=i)=p_i$.
    \STATE Compute partial gradient $\displaystyle g_t\gets \nabla_{i_t} f(w_t)$.
    \STATE Update the selected block:
    \[
    w_{t+1}^{(i_t)} \gets w_{t}^{(i_t)} - \eta\, g_t .
    \]
    \STATE Keep all other blocks unchanged.
\ENDFOR
\STATE \textbf{Return} $w_T$ (or optionally the average iterate $\bar w_T=\frac1T\sum_{t=0}^{T-1}w_t$).
\end{algorithmic}
\end{algorithm}

\section*{Dry Run Example}
Consider the one‑dimensional quadratic $f(w)=(w-3)^2$.
Here $B=1$ (single block), $\nabla f(w)=2(w-3)$, choose $\eta=0.1$, $w_0=0$.

\begin{center}
\begin{tabular}{c|c|c|c}
$t$ & $w_t$ & $\nabla f(w_t)$ & $f(w_t)$\\ \hline
0 & $0$ & $2(0-3)=-6$ & $9$\\
1 & $0-0.1(-6)=0.6$ & $2(0.6-3)=-4.8$ & $(0.6-3)^2=5.76$\\
2 & $0.6-0.1(-4.8)=1.08$ & $2(1.08-3)=-3.84$ & $(1.08-3)^2=3.6864$\\
3 & $1.08-0.1(-3.84)=1.464$ & $2(1.464-3)=-3.072$ & $(1.464-3)^2=2.3685\ (\text{approx})$
\end{tabular}
\end{center}
The objective value monotonically decreases, illustrating (1).

\newpage
\section*{Assumptions}
\begin{assumption}[Block‑wise $L_i$‑smoothness]\label{ass:smooth}
For every block $i\in\{1,\dots ,B\}$ there exists $L_i>0$ such that for any $w\in\mathbb{R}^d$ and any $u\in\mathbb{R}^{d_i}$,
\[
f\bigl(w^{(1)},\dots ,w^{(i)}+u,\dots ,w^{(B)}\bigr)
\le
f(w) + \langle \nabla_i f(w),u\rangle + \frac{L_i}{2}\|u\|^2 .
\tag{2}
\]
\end{assumption}
Define $L_{\max}\triangleq\max_{i}L_i$.

\begin{assumption}[Bounded Below]\label{ass:lower}
The function $f$ is bounded below: there exists $f_{\inf}>-\infty$ such that $f(w)\ge f_{\inf}$ for all $w$.
\end{assumption}

\begin{assumption}[Stepsize]\label{ass:stepsize}
The stepsize is constant $\eta_t\equiv\eta$ satisfying
\[
0<\eta\le \frac{1}{L_{\max}} .
\tag{3}
\]
\end{assumption}

\section*{Main Convergence Theorem}
\begin{theorem}[Expected Gradient Norm Decay]\label{thm:main}
Let $\{w_t\}_{t\ge0}$ be generated by Algorithm~\ref{alg:RBCD} under Assumptions~\ref{ass:smooth}--\ref{ass:stepsize}.  
Define the full gradient norm
\[
\|\nabla f(w_t)\|^2\;=\;\sum_{i=1}^{B}\|\nabla_i f(w_t)\|^2 .
\]
Then for any $T\ge1$,
\[
\frac{1}{T}\sum_{t=0}^{T-1}\mathbb{E}\!\bigl[\|\nabla f(w_t)\|^2\bigr]
\;\le\;
\frac{2\,B\bigl(f(w_0)-f_{\inf}\bigr)}{\eta\,T}.
\tag{4}
\]
Consequently,
\[
\min_{0\le t<T}\mathbb{E}\!\bigl[\|\nabla f(w_t)\|^2\bigr]\;=\;O\!\left(\frac{1}{T}\right).
\]
\end{theorem}

\section*{Theoretical Properties}
\begin{enumerate}[label=\arabic*.]
\item \textbf{Stability.} Because $\eta\le1/L_{\max}$, each update (1) is a descent step for the selected block (by (2)). Hence the iterates remain in a level set $\{w\mid f(w)\le f(w_0)\}$, which is bounded under mild coercivity assumptions.

\item \textbf{Hyper‑parameter Sensitivity.} The bound (4) scales linearly with the number of blocks $B$ and inversely with the stepsize $\eta$. Choosing $\eta$ close to $1/L_{\max}$ yields the fastest guaranteed $O(1/T)$ rate.

\item \textbf{Comparison to Full Gradient Descent.} For $B=1$ the bound reduces to the classic non‑convex GD rate 
$\frac{2\bigl(f(w_0)-f_{\inf}\bigr)}{\eta T}$.  
When $B$ grows, the factor $B$ appears because only one block is updated per iteration; the same total computational budget can be achieved by $B$ times more iterations.

\item \textbf{Special Cases.}
\begin{itemize}
\item \emph{Uniform sampling} ($p_i=1/B$) yields the clean constant $B$ in (4).
\item \emph{Non‑uniform sampling} with probabilities $p_i\propto L_i$ leads to a refined bound where $B$ is replaced by $\sum_i 1/p_i$, see Lemma~\ref{lem:nonuniform} below.
\end{itemize}
\end{enumerate}

\section*{Discussion}
The theorem shows that even without convexity, random block coordinate descent attains the same $O(1/T)$ decay of the expected squared gradient norm as full‑gradient stochastic methods, provided the stepsize respects the blockwise smoothness constants. The result is tight in the sense that the dependence on $B$ cannot be improved without additional structure (e.g., strong convexity or variance reduction).

\newpage
\section*{Preliminaries (Definitions and Lemmas)}
\begin{lemma}[Descent Lemma for a Single Block]\label{lem:descent}
Under Assumption~\ref{ass:smooth}, for any $w\in\mathbb{R}^d$ and any block $i$,
\[
f\!\bigl(w - \eta e_i \nabla_i f(w)\bigr)
\le
f(w) - \eta \|\nabla_i f(w)\|^2 + \frac{L_i\eta^{2}}{2}\|\nabla_i f(w)\|^2,
\tag{5}
\]
where $e_i$ denotes the embedding that updates only block $i$.
\end{lemma}
\begin{proof}
Apply (2) with $u=-\eta\nabla_i f(w)$:
\[
f\bigl(w^{(i)}-\eta\nabla_i f(w)\bigr)
\le
f(w)+\langle \nabla_i f(w), -\eta\nabla_i f(w)\rangle
+\frac{L_i}{2}\|\eta\nabla_i f(w)\|^{2}.
\]
Simplifying the inner product and the norm yields (5). \hfill $\square$
\end{proof}

\begin{lemma}[Expected One‑Step Decrease]\label{lem:expdecrease}
Let $i_t$ be sampled with probabilities $p_i$. Under (3) we have
\[
\mathbb{E}\!\bigl[f(w_{t+1})\mid w_t\bigr]
\le
f(w_t)-\frac{\eta}{2}\sum_{i=1}^{B}p_i\|\nabla_i f(w_t)\|^{2}.
\tag{6}
\]
\end{lemma}
\begin{proof}
Condition on $w_t$ and use Lemma~\ref{lem:descent}:
\[
\mathbb{E}\!\bigl[f(w_{t+1})\mid w_t\bigr]
= \sum_{i=1}^{B}p_i\,f\!\bigl(w_t-\eta e_i\nabla_i f(w_t)\bigr)
\le
\sum_{i=1}^{B}p_i\Bigl(f(w_t)-\eta\|\nabla_i f(w_t)\|^{2}
+\frac{L_i\eta^{2}}{2}\|\nabla_i f(w_t)\|^{2}\Bigr).
\]
Since $\eta\le 1/L_{\max}\le 1/L_i$ for every $i$, we have $L_i\eta^{2}\le\eta$. Hence
\[
-\eta+\frac{L_i\eta^{2}}{2}\le -\frac{\eta}{2}.
\]
Plugging this inequality gives (6). \hfill $\square$
\end{proof}

\begin{lemma}[Non‑uniform Sampling bound]\label{lem:nonuniform}
If $p_i>0$ and $\eta\le 1/L_{\max}$, then
\[
\mathbb{E}\!\bigl[f(w_{t+1})\mid w_t\bigr]
\le
f(w_t)-\frac{\eta}{2}\Bigl(\sum_{i=1}^{B}\frac{1}{p_i}\Bigr)^{-1}\|\nabla f(w_t)\|^{2}.
\tag{7}
\]
\end{lemma}
\begin{proof}
From Lemma~\ref{lem:expdecrease},
\[
\mathbb{E}[f(w_{t+1})\mid w_t]
\le f(w_t)-\frac{\eta}{2}\sum_{i=1}^{B}p_i\|\nabla_i f(w_t)\|^{2}.
\]
Apply Cauchy–Schwarz inequality in the form
\[
\Bigl(\sum_{i=1}^{B}p_i\|\nabla_i f\|^{2}\Bigr)
\ge
\Bigl(\sum_{i=1}^{B}\frac{1}{p_i}\Bigr)^{-1}
\Bigl(\sum_{i=1}^{B}\|\nabla_i f\|^{2}\Bigr)
= \Bigl(\sum_{i=1}^{B}\frac{1}{p_i}\Bigr)^{-1}\|\nabla f\|^{2}.
\]
Substituting yields (7). \hfill $\square$
\end{proof}

\section*{Main Proof (Step‑by‑Step Derivation)}
We now prove Theorem~\ref{thm:main}. All expectations are taken over the randomness of the selected blocks $\{i_t\}_{t=0}^{T-1}$.

\begin{proof}
\textbf{[Step 1]} Apply Lemma~\ref{lem:expdecrease} with uniform probabilities $p_i=1/B$:
\[
\mathbb{E}\!\bigl[f(w_{t+1})\mid w_t\bigr]
\le
f(w_t)-\frac{\eta}{2B}\sum_{i=1}^{B}\|\nabla_i f(w_t)\|^{2}
= f(w_t)-\frac{\eta}{2B}\|\nabla f(w_t)\|^{2}.
\tag{8}
\]

\textbf{[Step 2]} Take total expectation on both sides of (8):
\[
\mathbb{E}\!\bigl[f(w_{t+1})\bigr]
\le
\mathbb{E}\!\bigl[f(w_{t})\bigr]-\frac{\eta}{2B}\mathbb{E}\!\bigl[\|\nabla f(w_t)\|^{2}\bigr].
\tag{9}
\]

\textbf{[Step 3]} Rearrange (9) to isolate the gradient term:
\[
\frac{\eta}{2B}\mathbb{E}\!\bigl[\|\nabla f(w_t)\|^{2}\bigr]
\le
\mathbb{E}\!\bigl[f(w_{t})\bigr]-\mathbb{E}\!\bigl[f(w_{t+1})\bigr].
\tag{10}
\]

\textbf{[Step 4]} Sum inequality (10) over $t=0,\dots ,T-1$:
\[
\frac{\eta}{2B}\sum_{t=0}^{T-1}\mathbb{E}\!\bigl[\|\nabla f(w_t)\|^{2}\bigr]
\le
\sum_{t=0}^{T-1}\bigl(\mathbb{E}[f(w_{t})]-\mathbb{E}[f(w_{t+1})]\bigr)
= \mathbb{E}[f(w_{0})]-\mathbb{E}[f(w_{T})].
\tag{11}
\]

\textbf{[Step 5]} By Assumption~\ref{ass:lower}, $f(w_T)\ge f_{\inf}$. Hence
\[
\mathbb{E}[f(w_{0})]-\mathbb{E}[f(w_{T})]\le f(w_{0})-f_{\inf}.
\tag{12}
\]

\textbf{[Step 6]} Combine (11) and (12):
\[
\frac{\eta}{2B}\sum_{t=0}^{T-1}\mathbb{E}\!\bigl[\|\nabla f(w_t)\|^{2}\bigr]
\le
f(w_{0})-f_{\inf}.
\tag{13}
\]

\textbf{[Step 7]} Divide both sides of (13) by $T$ and multiply by $2B/\eta$:
\[
\frac{1}{T}\sum_{t=0}^{T-1}\mathbb{E}\!\bigl[\|\nabla f(w_t)\|^{2}\bigr]
\le
\frac{2B\bigl(f(w_{0})-f_{\inf}\bigr)}{\eta\,T}.
\tag{14}
\]

Equation (14) is precisely the bound (4) claimed in the theorem. \hfill $\square$
\end{proof}

\section*{Rate Derivation (With Constants)}
From (14) we obtain an explicit constant:
\[
C\;=\;\frac{2B\bigl(f(w_{0})-f_{\inf}\bigr)}{\eta},
\qquad
\frac{1}{T}\sum_{t=0}^{T-1}\mathbb{E}\!\bigl[\|\nabla f(w_t)\|^{2}\bigr]\le\frac{C}{T}.
\]
If the stepsize is taken at the maximal admissible value $\eta=1/L_{\max}$,
\[
C = 2B L_{\max}\bigl(f(w_{0})-f_{\inf}\bigr).
\]

\section*{Special Cases}
\begin{enumerate}[label=\alph*.]
\item \textbf{Uniform sampling, constant stepsize.} The bound (4) holds directly.

\item \textbf{Non‑uniform sampling.} Using Lemma~\ref{lem:nonuniform} we obtain
\[
\frac{1}{T}\sum_{t=0}^{T-1}\mathbb{E}\!\bigl[\|\nabla f(w_t)\|^{2}\bigr]
\le
\frac{2\bigl(\sum_{i=1}^{B}1/p_i\bigr)\bigl(f(w_{0})-f_{\inf}\bigr)}{\eta\,T},
\]
which reduces to (4) when $p_i=1/B$.

\item \textbf{Strongly smooth with $L_i\equiv L$ and $\eta=1/L$.} The descent factor becomes exactly $1/(2B)$, giving the tightest possible $O(1/T)$ rate within this analysis.

\end{enumerate}

\end{document}